\pagestyle{fancy}

\section{Introduction}

\subsection{Problem and Motivation}
\label{sec:problem-motivation}

Contemporary large language models reach the user not as the products of pre-training alone but as the products of an extensive post-training pipeline that layers supervised instruction tuning, preference learning from human feedback, constitutional AI, and direct preference optimisation on top of a base model that has merely learned to predict next tokens on web-scale text \parencite{christiano_deep_2023, ouyang_training_2022, ziegler_fine-tuning_2020, bai_constitutional_2022, rafailov_direct_2024}. It is this post-training stack, not the base model, that gives the deployed system its characteristic behaviours: following instructions, answering questions in a particular register, refusing harmful requests, and staying in the role of an assistant. Yet most mechanistic interpretability research on large language models has concentrated on base models---from \textcite{olah_zoom_2020}'s Zoom In essay on circuits to the indirect object identification circuit of \textcite{wang_interpretability_2022}, the factual association edits of \textcite{meng_locating_2023}, and the feed-forward key--value memories of \textcite{geva_transformer_2021}---leaving the internal mechanics of post-training comparatively under-studied. The result is a gap between the layer of the language model that interpretability understands best and the layer of the language model that users actually interact with.

Closing this gap is not a matter of scientific completeness alone; it has concrete implications for alignment practice. If we do not understand at a representational level what instruction tuning does to a model, we cannot know whether an observed safety behaviour in a deployed model is a robust property of the model's computation or a brittle veneer that can be undone by further fine-tuning, jailbroken by adversarial prompting, or silently circumvented by deceptive alignment \parencite{zou_universal_2023, hubinger_sleeper_2024}. The current answer to these questions is largely behavioural: we probe the deployed model, observe whether it refuses harmful prompts, and draw conclusions. This is informative but it is an outside-the-box view. What we lack is an inside-the-box account of which internal representations the post-training stack builds, reuses, or rearranges, and which of them carry the weight of the behaviours we care about.

The mechanism of instruction tuning in particular is the subject of a specific, contested, and empirically answerable debate. One widely cited position is the superficial alignment hypothesis associated with \textcite{zhou_lima_2023}: the observation that a surprisingly small amount of high-quality instruction-following data is sufficient to produce strong instruction-following behaviour in a large base model, implying that instruction tuning may not add new representational content so much as teach the model to expose existing content in a task-appropriate register. A competing position, less often articulated as a named hypothesis but implicit in much of the alignment literature, is that instruction tuning installs new computational structure---new features, new circuits, or new decision procedures---that did not previously exist in the base model and that are necessary for the characteristic behaviours of the instruction-tuned variant. These two positions predict observably different things about the internal representations of the instruction-tuned model, and whether a specific feature in the instruction-tuned model's layer 13 residual stream has a counterpart in the base model's layer 13 residual stream is exactly the kind of observation that can distinguish them.

Until recently, testing this distinction at the representational level was not practical. The tools available to mechanistic interpretability---probing classifiers, activation patching, attention pattern inspection, residual stream decomposition---either measured aggregate behavioural properties of the model or studied individual circuits one at a time, and neither approach gave a vocabulary in which two models' entire representational states could be systematically compared. The arrival of Sparse Autoencoders as a practical interpretability tool, beginning with the toy-model superposition work of \textcite{elhage_toy_2022} and early monosemanticity results \parencite{cunningham_sparse_2023}, and maturing through the Gemma Scope release and its successors \parencite{lieberum_gemma_2024}, changes this picture. SAEs give a feature-level dictionary of a model's residual-stream activations, with each dictionary entry corresponding to a direction in activation space that tends to activate on semantically legible patterns. Two models in the same architectural family, or---crucially for the present work---two checkpoints of the same model before and after post-training, can now be equipped with SAEs trained on the same residual-stream location, and their feature dictionaries can be compared directly. A feature that exists in both dictionaries, points in the same direction, and activates on the same inputs is a feature that has been preserved across post-training. A feature that exists in only one of the two dictionaries, or whose direction or functional behaviour has shifted, is a feature that has been modified, added, or removed. This comparison is the specific operationalisation of the debate between the superficial-alignment reading and the new-capability reading of instruction tuning, and it is what the present thesis undertakes.

The motivation for studying instruction tuning specifically, rather than the full post-training stack, is both practical and scientific. Practically, instruction tuning is the first and simplest stage of post-training, and its effects can be isolated from those of reinforcement-learning-based stages provided that matched Sparse Autoencoders exist for both the base model and the supervised-fine-tuned checkpoint. The Gemma Scope 2 release is, at the time of writing, the only such matched-SAE release for any production-class language model, and the Gemma 3 1B base and instruction-tuned checkpoints are its only covered pair---a constraint that determines the model selection of this thesis and is discussed in detail in Subsection~\ref{subsec:models-saes}. Scientifically, instruction tuning is also the stage of post-training whose mechanism is most directly addressed by the superficial-alignment literature and by the recent line of mechanistic fine-tuning work \parencite{jain_mechanistically_2024, prakash_fine-tuning_2024, du_how_2025}, which means that a feature-level study of instruction tuning connects immediately to a live debate in the interpretability and alignment communities rather than starting one from scratch.

The research question that follows from this motivation is stated in its final form in Subsection~\ref{sec:prior-art} after the literature review, but can be previewed here: at the feature level, does instruction tuning introduce new features into a language model's internal representations, or does it reorganise existing ones? The rest of this thesis sets out to answer that question for one specific model pair, using one specific interpretability tool, at one specific layer, with three behaviourally grounded probe datasets, and with explicit attention to the robustness of the answer across reasonable perturbations of its free parameters. What the thesis finds, in anticipation of Section~\ref{sec:results}, is a descriptive picture consistent with the reorganisation reading: the behaviourally relevant features of the instruction-tuned model are predominantly reshaped versions of features already present in the base model, not new features introduced from scratch. The conclusion is narrower than the full debate demands---it speaks to one layer of one model and is correlational rather than causal---but it is a concrete, reproducible data point on a question that has so far been addressed mostly at the behavioural and circuit level, and it provides a methodology that can be reused as matched Sparse Autoencoders become available for further model pairs.

\subsection{Prior Art and Research Question}
\label{sec:prior-art}

This subsection reviews the prior work that frames the research question of this thesis. It is organised thematically rather than chronologically, in four subsubsections. Subsubsection~\ref{subsubsec:repr-foundations} covers the representational foundations on which this work rests: the linear representation hypothesis, superposition, and dictionary learning as responses to the problem that individual neurons are not the right unit of analysis for understanding what a language model has learned. Subsubsection~\ref{subsubsec:saes-tool} turns to Sparse Autoencoders as a concrete interpretability tool, tracing the line from early toy models through to the Gemma Scope releases that make the present work possible. Subsubsection~\ref{subsubsec:post-training-effects} surveys what interpretability currently knows about the effect of post-training on internal representations, dividing the landscape into the superficial-alignment line of argument, the mechanistic fine-tuning line of work, and the representation-engineering and activation-steering work that sits adjacent to both. Subsubsection~\ref{subsubsec:gap-rq} identifies the specific gap in this literature that the present thesis addresses and states the research question in its final form.

\subsubsection{Representational Foundations: Superposition, Linear Features, and the Case for Dictionary Learning}
\label{subsubsec:repr-foundations}

The unit-of-analysis problem in mechanistic interpretability is older than interpretability as a field: it was already implicit in early work on distributed word representations \parencite{mikolov_efficient_2013}, which showed that directions in an embedding space could carry semantically meaningful structure that could not be read off any individual coordinate. Individual neurons, or individual coordinates, are generally not the right unit of analysis because the same neuron participates in many unrelated computations. \textcite{radford_learning_2017} observed this in the context of character-level language models, where a single ``sentiment neuron'' was identified but turned out to be an exception rather than the rule. \textcite{olah_feature_2017, olah_zoom_2020} articulated the more general view that the right unit of analysis is the feature---a direction in activation space that corresponds to a semantically legible pattern---and that features need not correspond to individual neurons.

The formal explanation for why individual neurons often fail as a unit of analysis is the phenomenon of \emph{superposition}, characterised in depth by \textcite{elhage_toy_2022}. The argument proceeds from a simple observation: a neural network with $n$ neurons cannot in general represent more than $n$ orthogonal directions exactly, but it can represent many more than $n$ near-orthogonal directions approximately, provided the represented features are sparse, that is, inactive on most inputs. A sparse set of features can be packed into a lower-dimensional space with interference that is tolerable in expectation, and the resulting representation assigns each feature a direction that is in general a linear combination of several neurons. The consequence is that single-neuron interpretation is not merely aesthetically limited but methodologically incorrect for sufficiently overcomplete learned representations: the neurons of such a network will be polysemantic by design. \textcite{scherlis_polysemanticity_2025} extend this analysis to a more rigorous account of polysemanticity as a function of capacity, and \textcite{park_linear_2024} formalise the linear representation hypothesis---that meaningful features in language model activations correspond to directions rather than to nonlinear manifolds---as a testable structural claim about transformer internals.

Superposition motivates dictionary learning as a natural remedy. If the true representational units of the network are features that live as sparse linear combinations of neurons, then the correct way to recover them is to learn a dictionary of feature directions and an encoder that assigns each activation a sparse code over that dictionary. This is a classical idea in signal processing, and its application to neural network interpretability received its first major language-model demonstration in \textcite{cunningham_sparse_2023}, who showed that sparsely activated dictionaries trained on transformer activations recovered features substantially more interpretable than individual neurons on the same models. Concurrent work in Anthropic's Towards Monosemanticity line \parencite{bricken_monosemanticity_2023} scaled the approach to larger transformers and established much of the vocabulary---``features,'' ``monosemantic directions,'' and ``activating examples''---that has become standard in the SAE literature. The methodological premise of the present thesis rests on this line of work: that features recovered by dictionary learning on a transformer's residual stream are a meaningful representational unit, and that two transformers whose residual streams are dictionary-decoded in the same way can be compared at the level of their features.

\subsubsection{Sparse Autoencoders as an Interpretability Tool}
\label{subsubsec:saes-tool}

Sparse Autoencoders (SAEs) are the specific form of dictionary learning used in this thesis and in most recent SAE-based interpretability work. An SAE takes an activation $h \in \mathbb{R}^{D}$ from some location in a transformer, typically the residual stream at a specified layer, encodes it into a sparse non-negative feature vector $f(h) \in \mathbb{R}^{F}$ with $F > D$, and decodes it through a learned dictionary $W_{\mathrm{dec}} \in \mathbb{R}^{F \times D}$ to produce a reconstruction $\hat{h} = W_{\mathrm{dec}}^{\top} f(h)$. Training minimises reconstruction loss together with an explicit or implicit sparsity penalty, so that only a small subset of the $F$ features activate on any given input. The resulting decoder rows are the feature directions, and the encoder gives the sparse code.

Early SAE work used the classical formulation with an $L_1$ sparsity penalty on the encoder output, but this formulation has known failure modes. $L_1$ penalties shrink active features toward zero, biasing both reconstruction and downstream analysis, and the fixed sparsity regime is difficult to tune across layers and scales. Two important methodological refinements emerged in response. \textcite{rajamanoharan_improving_2024} introduce gated SAEs, which decouple the decision of whether a feature is active from the magnitude of its activation, reducing the shrinkage bias of the $L_1$ formulation, and later work in the same line introduces the JumpReLU encoder used in Gemma Scope, which makes the feature activation threshold explicit and trainable. \textcite{gao_scaling_2024} provide a systematic study of how SAE quality scales with dictionary width, training compute, and sparsity target on large language models, and establish empirical scaling laws that guided the hyperparameter choices of subsequent SAE releases. \textcite{karvonen_saebench_2025} propose a comprehensive evaluation benchmark, SAEBench, that measures SAE quality along multiple axes beyond reconstruction loss, including feature interpretability, probe-based downstream utility, and robustness to perturbation.

The most important concrete artefact for the present thesis is the Gemma Scope release \parencite{lieberum_gemma_2024}. Gemma Scope made publicly available a large family of trained Sparse Autoencoders across every layer and several widths of Gemma-family models, and Gemma Scope 2 extended this to include matched releases for the Gemma 3 1B base and instruction-tuned checkpoints used throughout this thesis. The scale and reproducibility of Gemma Scope changed what was practical to ask of SAE-based interpretability: researchers could now write papers that assumed the existence of high-quality trained SAEs for a production-class language model rather than training their own from scratch. \textcite{peng_use_2025} further argues that SAEs are most useful for discovering unknown concepts rather than acting on previously known ones, a caution that is relevant to the targeted analysis stage of the present thesis, which uses behavioural probe datasets to let the feature space surface its own concepts rather than prescribing them in advance.

\subsubsection{Post-Training and Its Effects on Internal Representations}
\label{subsubsec:post-training-effects}

Post-training is the umbrella term for everything that happens to a language model after pre-training on next-token prediction. The standard modern pipeline begins with supervised instruction tuning \parencite{wei_finetuned_2022, zhou_lima_2023}, continues with reward modelling \parencite{christiano_deep_2023, ziegler_fine-tuning_2020}, and often concludes with reinforcement learning from human feedback \parencite{ouyang_training_2022}. Constitutional AI \parencite{bai_constitutional_2022} and related variants replace or augment the human feedback with model-generated feedback grounded in an explicit set of principles, and \textcite{rafailov_direct_2024}'s Direct Preference Optimisation reformulates the RL stage as a direct optimisation problem that sidesteps the reward model altogether. The broader helpful-and-harmless line of \textcite{bai_training_2022} frames the goal of this pipeline as producing a model that is simultaneously useful to users and safe to deploy.

The question of what all this does to the internal representations of the model has been addressed from three distinct angles in the interpretability literature, and the research question of this thesis sits at the intersection of all three.

\textbf{The superficial alignment line of argument.} The strongest statement of this view is \textcite{zhou_lima_2023}'s LIMA paper, whose central empirical finding is that 1,000 carefully curated instruction-response examples are sufficient to produce strong instruction-following behaviour in a large base model, comparable to models fine-tuned on orders of magnitude more data. The authors' interpretation is explicit: a model's knowledge and capabilities are learned almost entirely during pre-training, while alignment teaches it which subdistribution of formats should be used when interacting with users. This is a behavioural-evidence argument and does not itself commit to any particular representational mechanism, but it makes a strong prediction about what an interpretability study would find: that the substrate for instruction-following behaviour is already present in the base model and that fine-tuning primarily affects which features are read out and in what register, rather than which features exist. The LIMA paper does not test this prediction at the representational level, but it frames the exact empirical question that the present thesis answers.

\textbf{The mechanistic fine-tuning line.} A parallel line of work approaches the same question from a mechanistic rather than behavioural angle, using circuit-level and probe-level tools to ask what structures inside the model change when it is fine-tuned. The foundational paper for the present thesis is \textcite{jain_mechanistically_2024}, which studies fine-tuning in controlled synthetic settings using compiled Tracr transformers and probabilistic context-free grammars, and finds that fine-tuning rarely alters underlying model capabilities, instead learning a minimal transformation, or ``wrapper,'' on top of them. The wrapper hypothesis---that fine-tuning installs a minimal re-routing of existing capabilities rather than building new ones---is the specific claim this thesis tests in a new setting. Jain et al.'s evidence comes from pruning, probing, and reverse fine-tuning in synthetic environments where ground truth is known; it does not come from SAE-based feature-level analysis and does not apply directly to a production-class language model. Extending the wrapper hypothesis to a real instruction-tuned LLM at the feature level is the specific extension attempted here.

Adjacent to Jain et al. is \textcite{prakash_fine-tuning_2024}, which uses circuit-level tools in a more realistic setting to show that fine-tuning for entity tracking does not build a new circuit but rather enhances the selectivity and magnitude of a circuit already present in the base model. \textcite{du_how_2025} extends this picture further and is directly concerned with which aspects of model behaviour are most affected by post-training at the representational level. \textcite{xu_tracking_nodate} works at a different temporal resolution, within pre-training itself, but contributes the methodological observation that feature-level analysis can track representational change over training, which is the temporal analogue of the checkpoint-comparison method used in this thesis.

\textbf{Representation engineering, activation steering, and the causal side.} A third adjacent literature uses representations not only as objects of interpretation but as targets of intervention. \textcite{zou_representation_2025} propose a top-down approach in which high-level concept directions are extracted from model activations using contrast pairs and then used to monitor, probe, and steer model behaviour, demonstrating that properties such as honesty, refusal, and emotional valence can be read from and written to the residual stream via linear operations. This line of work provides the causal counterpart to correlational feature-level analyses: where SAE-based work asks what features exist in a model, representation engineering asks what happens when those features are perturbed. The present thesis makes a limited excursion into this territory via the steering analysis reported in Subsection~\ref{subsec:rq3-results}.

The jailbreak and adversarial-attack literature \parencite{zou_universal_2023, hubinger_sleeper_2024} provides the motivating edge for this line of work. If safety behaviours are a thin wrapper over a base-model capability substrate, then adversarial prompts that bypass the wrapper should be able to surface capabilities that the instruction-tuned model would otherwise not exhibit. A feature-level picture of which representations are shared between base and instruction-tuned checkpoints, which is what this thesis provides, is directly relevant to understanding where this surfacing can and cannot reach.

\textbf{Similarity-of-representations work.} One further piece of adjacent literature deserves mention: \textcite{kornblith_similarity_2019}'s work on Centered Kernel Alignment (CKA) and related similarity metrics for comparing neural network representations. CKA and its variants give an aggregate measure of how similar two networks' representations are at a given layer, and have been used extensively to compare fine-tuned models to their base models. Aggregate similarity measures are complementary to the feature-level comparison of this thesis: CKA tells you how much two representations differ overall, whereas a feature-level analysis tells you which specific axes differ and how. The present thesis does not compute CKA between the PT and IT residual streams, but the feature-level results reported later could in principle be sanity-checked against such a computation in future replications.

\subsubsection{The Gap and the Research Question}
\label{subsubsec:gap-rq}

The literature reviewed in Subsubsections~\ref{subsubsec:repr-foundations}--\ref{subsubsec:post-training-effects} establishes three things. First, Sparse Autoencoders are now a sufficiently mature tool to give a feature-level vocabulary for the residual stream of a production-class language model, and matched SAEs exist for the Gemma 3 1B base and instruction-tuned checkpoints via Gemma Scope 2. Second, there is a live and contested hypothesis that instruction tuning primarily re-exposes existing base-model capabilities rather than installing new ones, most prominently in the LIMA paper but implicit across much of the alignment literature. Third, this hypothesis has already been tested with non-SAE tools in synthetic settings and with circuit-level tools in specific task settings, and the results of both lines of work are broadly consistent with it in their respective domains, but it has not been tested at the feature level in a real instruction-tuned LLM using matched Sparse Autoencoders.

The specific gap this thesis addresses is the intersection of these observations: there is no existing work that uses matched Sparse Autoencoders to compare the feature-level representations of a real pre-trained and instruction-tuned language model at the same layer and to test, descriptively, whether the behaviourally relevant features of the instruction-tuned model are preserved, modified, or newly introduced relative to the base model. The Gemma Scope 2 release makes this comparison possible for the first time on Gemma 3 1B, and the absence of prior work in this intersection is what motivates the pipeline described in Section~\ref{sec:methodology}.

The research question that follows is: at the feature level, does instruction tuning introduce new features into a language model's internal representations, or does it predominantly reorganise existing ones? Specifically, when the residual stream of a pre-trained and an instruction-tuned variant of Gemma 3 1B is decoded with matched Sparse Autoencoders at layer 13, are the features most behaviourally relevant to instruction-following and safety-related prompts better characterised as preserved, modified, or exclusive to one of the two checkpoints?

The question is deliberately narrower than the full superficial-alignment versus new-capability debate. It restricts itself to a single layer, a single model pair, and a single post-training stage. It also restricts itself primarily to a descriptive and correlational framing, even though the thesis includes a limited causal steering check, for reasons of scope discussed in Section~\ref{sec:methodology} and Subsection~\ref{sec:limitations-robustness}. The narrowness is intentional: a question that can be answered precisely with the available tools is more valuable than a broader question that cannot. The answer to this narrower question, as later sections report and interpret, provides one concrete feature-level data point in a debate that has so far been conducted mostly at the behavioural and circuit level.

This question is operationalised through three sub-questions:
\begin{itemize}
    \item RQ1: What proportion of features at layer 13 of Gemma 3 1B are preserved, modified, PT-exclusive, or IT-exclusive between the base and IT checkpoints, and how robust is that proportion to the choice of matching thresholds?
    \item RQ2: When the IT model is probed with safety-relevant and instruction-following prompts, do the most differentially active features map onto IT-exclusive entries, as a strict-creation reading would predict, or onto features that the matching pipeline already paired with a PT counterpart?
    \item RQ3: Does causal manipulation via activation steering of the features identified in RQ2 produce the behavioural changes that their labels predict, in both the base and IT models?
\end{itemize}

\subsection{Main Contribution}
\label{sec:main-contribution}

The main contribution of this thesis is a four-stage empirical pipeline for testing whether instruction tuning introduces new features into a language model's internal representations or predominantly reorganises pre-existing ones, together with the first application of this pipeline to a real instruction-tuned language model using matched Sparse Autoencoders. Applied to the pre-trained and instruction-tuned variants of Gemma 3 1B at the residual stream of layer 13, using the \texttt{gemma-scope-2-1b-pt-res} and \texttt{gemma-scope-2-1b-it-res} Sparse Autoencoders from the Gemma Scope 2 release, the pipeline produces empirical evidence that the features most behaviourally relevant to instruction-following and safety-related prompts are predominantly reorganised pre-existing features rather than newly introduced ones---a descriptive result consistent with the wrapper hypothesis of \textcite{jain_mechanistically_2024}, here tested for the first time in a real instruction-tuned LLM at the feature level.

The concrete deliverables of the thesis, each of which is elaborated in the chapters that follow, are as follows.

The first deliverable is the \textbf{four-stage pipeline itself}. Stage one collects residual-stream activations from 5,000 FineWeb sequences for both the pre-trained and the instruction-tuned checkpoints at layer 13 and encodes them with the matched Gemma Scope 2 Sparse Autoencoders, producing per-feature activation frequencies, mean activations, and per-sequence mean activation vectors for each model. Stage two computes two complementary feature-correspondence measures between the two checkpoints---decoder cosine similarity and activation correlation---and uses them jointly to partition the active feature sets of both checkpoints into four classes: \textit{preserved} features (strong correspondence on both measures), \textit{modified} features (strong correspondence on only one of the two measures), \textit{PT-exclusive} features (active in the base model with no strong counterpart in the instruction-tuned dictionary), and \textit{IT-exclusive} features (active in the instruction-tuned model with no strong counterpart in the base dictionary). The classification is subjected to a 25-point threshold sensitivity grid to verify that its qualitative shape is not an artefact of the primary operating point chosen for the interpretive stages. Stage three attaches a descriptive vocabulary to a frequency-weighted sample of 50 features from each class by submitting their top-activating contexts to two independent large language model annotators (Claude Haiku 4.5 and GPT-5.4-nano) under identical prompts and identical structured-output schemas, and constructs a consensus label by confidence-weighted tiebreak on disagreements, treating the inter-annotator agreement rate and the structure of disagreement cases as first-class reliability diagnostics rather than as noise to be filtered away. Stage four runs a targeted differential activation analysis on three external behavioural probe datasets---JailbreakBench for safety-relevant prompts, Stanford Alpaca for instruction-following prompts, and TruthfulQA for a truthfulness contrast---using Welch's two-sample $t$-tests with an explicit $p < 0.01$ threshold and a minimum conditional-mean activation filter, and cross-references the resulting ranked lists of differential features against the four-way classification of stage two.

The second deliverable is the empirical findings that the pipeline produces. These are reported in full in Chapter 6 and summarised in Chapter 8, and consist of three convergent strands of descriptive evidence. The classification of stage two finds that cleanly preserved features are a small minority of active features at the primary operating point, that the modified class is dominant across the entire threshold sensitivity grid, and that both checkpoint-exclusive classes are sizeable but---as the targeted analysis of stage four shows---not behaviourally central. The dual-LLM consensus labelling of stage three finds that inter-annotator agreement is moderate overall and is concentrated on the preserved and modified classes, that disagreement concentrates at specific identifiable taxonomic boundaries (most prominently the boundary between low-level linguistic features and instruction-style features), that preserved features are almost entirely low-level linguistic and syntactic in character, and that instruction-style features are essentially absent from the preserved class and concentrated in the modified class instead. The targeted differential analysis of stage four finds that the top-ranked features responding to instruction-shaped prompts in the instruction-tuned model are overwhelmingly labelled as instruction-style and that several of them---most saliently, specific feature indices that can be identified by number in Chapter 6---appear in the corresponding top-ranked lists for the base model on the same prompts with large effect sizes, meaning that these features are already present and already selective in the base model. The same cross-checkpoint overlap pattern holds, more weakly, for features responding to harmful prompts under the JailbreakBench contrast.

The third deliverable is the methodological template that the pipeline represents. The four stages are loosely coupled and can be reused independently: the feature-matching step of stage two can be applied to any pair of matched SAEs; the dual-LLM consensus labelling of stage three can be applied to any set of SAE features with top-activating contexts; the targeted differential analysis of stage four can be applied to any set of features identified by any other means. The template is the first published application of matched Sparse Autoencoder comparison between a base and an instruction-tuned checkpoint of a production-class language model, and it defines a reusable protocol for any future study that will have access to matched SAEs for a different model pair.

What the thesis does not claim is equally important and equally explicit. The descriptive evidence supports a correlational reading, not a causal one. Chapter 7 is explicit about the absence of a behaviourally grounded causal validation experiment, reports the partial results of a cross-SAE direction robustness check as a structural observation rather than as a causal finding, and names the absence of full causal validation as the primary limitation of the work and the most important piece of future work. Chapter 8's verdict on the wrapper hypothesis is therefore a calibrated ``consistent with'' rather than a confirmatory ``supports,'' and the strongest version of the hypothesis---that instruction tuning creates no new features at all---is explicitly rejected as unsupported by the evidence. What the thesis does claim, and what the pipeline and findings support, is the weaker and more defensible position that at the residual stream of layer 13 in Gemma 3 1B, the features carrying the behavioural signatures of instruction tuning are predominantly reorganised from a pre-existing base-model feature population rather than introduced from scratch. This is the contribution, and the rest of the thesis is the defence of it.
